\chapter{Vaginal atrophy atrophic vaginitis}
\index{Vaginal atrophy atrophic vaginitis}

\section*{Definition and Overview}
Vaginal atrophy, or atrophic vaginitis, is a chronic, progressive condition characterized by the thinning, drying, and inflammation of the vaginal walls due to declining estrogen levels. In modern medicine, it is classified under the Genitourinary Syndrome of Menopause (GSM), a term introduced in 2014 by the International Society for the Study of Women's Sexual Health and the Menopause Society. GSM better reflects the comprehensive anatomical, physiological, and functional changes that estrogen deficiency causes across the labia, clitoris, vagina, urethra, and bladder, moving away from the restrictive term ``atrophic vaginitis'' which implies only vaginal inflammation. Estrogen is vital for maintaining the elasticity, vascularity, and moisture of the urogenital tissues. It stimulates epithelial cell proliferation and glycogen production, which supports Lactobacilli to maintain an acidic vaginal pH (3.8 to 4.5). Estrogen depletion disrupts this equilibrium, causing thinning of the vaginal lining, loss of elasticity, and mucosal dryness. This state leads to significant physical discomfort, including dyspareunia (painful intercourse), vulvar burning, and urinary symptoms, profoundly impacting the quality of life, sexual health, and emotional well-being of affected individuals.

\section*{Epidemiology}
Vaginal atrophy is a highly prevalent condition that affects approximately 50\% to 60\% of postmenopausal women, though clinical findings suggest up to 90\% exhibit signs of tissue regression. Unlike vasomotor symptoms (hot flushes) which typically improve over time, GSM is progressive and worsens without intervention.
The primary affected populations include:
\begin{itemize}
\item **Natural Menopause**: Women undergoing the natural aging process, with symptoms usually appearing during perimenopause and accelerating postmenopausally.
\item **Surgical Menopause**: Women who undergo bilateral oophorectomy, experiencing an abrupt, severe drop in estrogen.
\item **Premature Ovarian Insufficiency**: Women experiencing ovarian failure before age 40.
\item **Oncology Patients**: Women undergoing chemotherapy, pelvic radiation, or endocrine therapy (such as aromatase inhibitors or tamoxifen) for breast cancer, which severely suppresses estrogen levels.
\item **Postpartum and Lactating Women**: Individuals experiencing temporary hypoestrogenism due to elevated prolactin levels.
\end{itemize}
Demographically, GSM affects individuals globally across all races and ethnicities, though socioeconomic and cultural barriers often lead to underreporting and delayed treatment.

\section*{Aetiology and Pathophysiology}
The primary cause of vaginal atrophy is the decline in circulating estrogen levels. The pathophysiology involves several key urogenital changes:
\begin{itemize}
\item **Epithelial Thinning**: Without estrogen, the stratified squamous epithelium of the vagina thins down to a fragile layer of parabasal cells, losing its protective superficial cells and mucosal barrier integrity.
\item **Loss of Glycogen and Lactobacilli**: Estrogen stimulates glycogen synthesis in vaginal cells. When glycogen shedding drops, the population of acid-producing Lactobacilli (such as \textit{Lactobacillus} species) declines.
\item **Elevation of Vaginal pH**: The loss of lactic acid causes vaginal pH to rise from an acidic 3.8 to 4.5 to an alkaline 5.5 to 7.5, facilitating colonization by opportunistic pathogens (e.g., \textit{Escherichia coli}, \textit{Streptococcus}) and increasing vaginitis risk.
\item **Reduced Vascularity and Elasticity**: Decreased pelvic blood flow reduces vaginal lubrication and secretions. Collagen and elastin fibers in the lamina propria degrade, causing tissue fibrosis, narrowing, and shortening of the vaginal canal.
\item **Lower Urinary Tract Changes**: The urethra and bladder trigone share a common embryological origin with the vagina and contain high levels of estrogen receptors. Estrogen depletion thins the urethral mucosa and decreases detrusor compliance, leading to urinary symptoms.
\end{itemize}

\section*{Clinical Features}
The clinical presentation of vaginal atrophy involves a range of vulvovaginal, sexual, and urinary signs and symptoms:
Symptoms include:
\begin{itemize}
\item **Vaginal Dryness**: Persistent dryness, itching, or burning in the vulvovaginal area.
\item **Dyspareunia**: Pain during sexual activity, ranging from mild discomfort to severe tearing pain, often leading to secondary vaginismus.
\item **Abnormal Discharge**: Thin, watery, yellow or grey discharge, occasionally blood-tinged, with an altered odour due to elevated pH.
\item **Urinary Symptoms**: Increased urinary frequency, urgency, nocturia, and dysuria.
\item **Recurrent Infections**: Frequent urinary tract infections due to loss of protective acidic flora.
\end{itemize}
Clinical signs observed during pelvic examination include:
\begin{itemize}
\item **Mucosal Pallor and Dryness**: The vaginal walls appear pale, shiny, and dry.
\item **Loss of Rugae**: Flattening or complete absence of the normal vaginal folds.
\item **Petechiae and Friability**: Pinpoint hemorrhages and tissues that bleed easily upon contact.
\item **Vaginal Stenosis**: Shortening and narrowing of the vaginal canal and constriction of the introitus.
\item **Urethral Caruncle**: A small, red, benign vascular growth at the posterior urethral meatus.
\end{itemize}

\section*{Differential Diagnosis}
A thorough differential diagnosis is essential to rule out other causes of vulvovaginal discomfort:
\begin{itemize}
\item **Infectious Vaginitis**: Candidiasis (typically characterized by a thick, white discharge and acidic pH), bacterial vaginosis (causing a thin, fishy-smelling discharge), and trichomoniasis.
\item **Lichen Sclerosus**: A chronic inflammatory dermatosis causing ivory-white plaques and structural changes on the vulva. Unlike GSM, it does not affect the vaginal canal itself.
\item **Erosive Lichen Planus**: An autoimmune condition causing painful mucosal erosions, desquamative discharge, and vaginal adhesions.
\item **Contact Dermatitis**: Irritation or allergy of the vulva caused by hygiene products, soaps, or synthetic fabrics.
\item **Vulvodynia**: Chronic, unexplained vulvar pain or burning lasting at least 3 months.
\item **Pelvic Floor Muscle Dysfunction**: Hypertonic pelvic floor muscles causing dyspareunia and pelvic pain, often secondary to chronic atrophic pain.
\item **Malignancy**: Vulvar or vaginal intraepithelial neoplasia, squamous cell carcinoma, or endometrial cancer (which must be suspected in any postmenopausal bleeding).
\end{itemize}

\section*{When to Seek Medical Care}
While vaginal atrophy is benign, patients must seek medical evaluation if they experience concerning symptoms. Immediate consultation is required for any postmenopausal bleeding, spotting, or unusual discharge occurring more than 12 months after the final menstrual period, to rule out endometrial neoplasia. Patients should also see a doctor for unexplained vulvar lesions, severe pelvic pain, or signs of urinary tract infection (dysuria, fever, hematuria).
For emergency symptoms, such as high fever with chills, severe back or flank pain (indicating pyelonephritis), or heavy vaginal bleeding, patients should seek immediate emergency care by dialing:
\begin{itemize}
\item **local emergency services** in many countries.
\item **911** in the United States and Canada.
\item **999** in the United Kingdom.
\item **112** in the European Union.
\end{itemize}

\section*{Diagnosis and Investigations}
The diagnosis of vaginal atrophy is established through a clinical history, physical exam, and specific investigations:
\begin{itemize}
\item **Physical and Speculum Exam**: Evaluates the appearance of the vulva and vagina for mucosal pallor, loss of rugae, petechiae, and stenosis.
\item **Vaginal pH Testing**: A pH strip is placed against the lateral vaginal wall; a pH greater than 5.0 supports the diagnosis in the absence of semen or infection.
\item **Vaginal Maturation Index**: A cytological smear of the vaginal wall. A high percentage of parabasal cells and low percentage of superficial cells confirms severe estrogen deficiency.
\item **Wet Mount Microscopy**: Excludes yeast, bacterial vaginosis, or trichomoniasis.
\item **Urinalysis and Urine Culture**: Performed to rule out active urinary tract infections in symptomatic patients.
\item **Transvaginal Ultrasound and Endometrial Biopsy**: Indicated if postmenopausal bleeding occurs, evaluating endometrial thickness and checking for endometrial pathology.
\end{itemize}

\section*{Management}
Treatment is tailored to symptom severity, patient history, and preferences:
\begin{itemize}
\item **Vaginal Moisturizers**: Non-hormonal, over-the-counter products applied 2 to 3 times weekly to restore mucosal hydration and tissue elasticity.
\item **Vaginal Lubricants**: Water- or silicone-based lubricants used on-demand during sexual intercourse to reduce friction.
\item **Local Estrogen Therapy**: The gold standard for moderate-to-severe symptoms. Very low doses applied locally restore the epithelium with minimal systemic absorption. Standard forms include:
\begin{itemize}
\item **Creams**: Estradiol or conjugated estrogens creams, allowing flexible dosing.
\item **Tablets**: Low-dose estradiol tablets inserted twice weekly after a loading phase.
\item **Rings**: A flexible silicone ring releasing estradiol consistently over 90 days.
\item **Inserts**: Ultra-low-dose estradiol softgel inserts.
\end{itemize}
\item **Systemic Hormone Replacement Therapy**: Indicated when systemic menopausal symptoms (e.g., hot flushes) coexist with vaginal atrophy.
\item **Oral Selective Estrogen Receptor Modulators**: **Ospemifene** is a daily oral tablet that acts as an estrogen agonist on the vaginal tissue, improving cell thickness and reducing dyspareunia.
\item **Intravaginal DHEA (Prasterone)**: A daily insert that undergoes local intracellular conversion into estrogens and androgens.
\item **Vaginal Laser Therapy**: Fractional carbon dioxide or Erbium:YAG lasers designed to stimulate collagen production and mucosal regeneration.
\item **Pelvic Floor Physical Therapy**: Helps desensitize tissues and relax hypertonic muscles associated with chronic painful intercourse.
\end{itemize}

\section*{Complications}
If left untreated, vaginal atrophy can lead to severe and chronic complications:
\begin{itemize}
\item **Vaginal and Introital Stenosis**: Narrowing and shortening of the vaginal canal, which can make intercourse and pelvic examinations impossible.
\item **Sexual Dysfunction**: Severe dyspareunia leading to secondary vaginismus, loss of libido, and relationship distress.
\item **Recurrent Urinary Infections**: Chronic ascending infections due to alkaline pH and urethral thinning, which can progress to pyelonephritis or urosepsis.
\item **Tissue Fissuring**: Fragile tissues prone to tearing, bleeding, and secondary infections.
\item **Incontinence Exacerbation**: Worsening of stress and urge urinary incontinence.
\item **Treatment Risks**: Local estrogen may cause mild vaginal discharge. Systemic hormone replacement therapy carries small risks of breast cancer, thromboembolism, or stroke depending on duration and formulation.
\end{itemize}

\section*{Prognosis and Outcomes}
The prognosis for vaginal atrophy is excellent when managed with appropriate, consistent treatment. Because GSM is a chronic and progressive condition, symptoms will recur if therapies are stopped.
Key prognoses include:
\begin{itemize}
\item **Response Time**: Vaginal dryness and burning typically improve within 2 to 4 weeks, while dyspareunia resolves within 8 to 12 weeks of initiating local estrogen.
\item **Urinary Recovery**: The frequency of urinary tract infections and symptoms of urgency decrease as the vaginal microbiome and urethral tissues heal.
\item **Maintenance**: Long-term adherence is required. Low-dose local estrogen can be safely used for years to maintain urogenital health.
\item **Early Intervention**: Early treatment prevents irreversible structural changes like severe vaginal stenosis.
\end{itemize}

\section*{Prevention and Self-Care}
Self-care measures can delay the onset and mitigate the severity of vaginal atrophy:
\begin{itemize}
\item **Regular Sexual Activity**: Engaging in regular sexual activity (including vaginal intercourse, masturbation with vaginal penetration, or the use of vaginal dilators) is highly effective in maintaining vaginal health. Sexual stimulation promotes blood flow to the pelvic organs, increases natural lubrication, and helps maintain the length, width, and elasticity of the vaginal canal.
\item **Avoidance of Chemical Irritants**: Refraining from using perfumed soaps, bubble baths, douches, feminine hygiene sprays, and scented laundry detergents.
\item **Hygiene Practices**: Washing the vulva with warm water only, and wiping from front to back to prevent bacterial contamination.
\item **Clothing Choices**: Wearing loose-fitting clothing and 100\% cotton underwear to avoid trapping moisture.
\item **Smoking Cessation**: Smoking accelerates estrogen metabolism and reduces pelvic blood flow; quitting improves tissue perfusion and reduces symptom severity.
\item **Hydration and Diet**: Maintaining adequate hydration and eating a balanced diet rich in essential fatty acids.
\end{itemize}

\section*{Sources and Further Reading}
For additional clinical guidelines and patient resources, consult:
\begin{itemize}
\item **The Menopause Society**: Provides resources on midlife health and GSM. URL: https://www.menopause.org
\item **American College of Obstetricians and Gynecologists**: Clinical guidelines on menopause management. URL: https://www.acog.org
\item **Australasian Menopause Society**: Fact sheets and clinical updates on vaginal health. URL: https://www.menopause.org.au
\item **British Menopause Society**: Guidelines and consensus statements on postmenopausal care. URL: https://www.thebms.org.uk
\item **International Society for the Study of Women's Sexual Health**: Research and guidelines on sexual dysfunction. URL: https://www.isswsh.org
\item **Mayo Clinic**: Patient-friendly medical overviews of atrophic vaginitis. URL: https://www.mayoclinic.org
\end{itemize}